\documentclass[12pt]{article}
\PassOptionsToPackage{table}{xcolor}
\usepackage[utf8]{inputenc}
\usepackage[T1]{fontenc}
\usepackage[italian]{babel}
\usepackage{multicol}
\usepackage{enumitem}
\usepackage{array}
\usepackage{longtable}
\usepackage{booktabs}
\usepackage{xcolor}
\usepackage{graphicx}
\usepackage{tikz}
\usetikzlibrary{
    shapes.geometric,
    arrows.meta,
    positioning,
    calc,
    patterns,
    decorations.pathmorphing
}
\usepackage[margin=2cm]{geometry}
\usepackage{amsmath,amssymb}
\usepackage{siunitx}
\setlength{\parindent}{0pt}
\definecolor{headercolor}{RGB}{220,220,220}
\newcounter{quesito}
\setcounter{quesito}{0}
\newcommand{\nuovoquesito}{\stepcounter{quesito}\textbf{Domanda \arabic{quesito}.}}
\newcommand{\itemdomanda}[2]{%
    \par\noindent
    \begin{minipage}[t]{\linewidth}
        \nuovoquesito \\ #1
        \begin{enumerate}[label=(\Alph*), nosep, leftmargin=*]
            #2
        \end{enumerate}
    \end{minipage}
    \vspace{10pt}
}
\begin{document}
\section*{Simulazione}
\hfill Data: 17 apr 2026

\itemdomanda{Se A è più alto di B e B è più alto di C, allora:}{
    \item C è più alto di A
    \item A è più basso di C
    \item A è più alto di C
    \item B è il più alto di tutti
    \item Non si può stabilire nulla
}

\itemdomanda{Date le seguenti premesse:
    \begin{enumerate}
        \item Tutti i canarini ben nutriti cantano a squarciagola.
        \item Nessun canarino che canti a squarciagola è malinconico.
    \end{enumerate}
    Quale conclusione è logicamente corretta?}{
    \item Nessun canarino ben nutrito è malinconico
    \item Qualche canarino ben nutrito è malinconico
    \item Tutti i canarini malinconici sono ben nutriti
    \item I canarini che non cantano sono sicuramente ben nutriti
    \item Chi canta a squarciagola è malinconico
}

\itemdomanda{Considerando i dati in tabella l'unica relazione possibile tra le grandezze $x$ e $y$ è di:
\renewcommand{\arraystretch}{1.5}
\begin{tabular}{|>{\columncolor{headercolor}}c|c|c|c|c|c|c|}
\hline
$x$ & 2 & 4 & 6 & 8 & 10 & 12 \\ \hline
$y$ & 1 & 4 & 7 & 10 & 13 & 16 \\ \hline
\end{tabular}

}{
\item proporzionalità quadratica
\item proporzionalità inversa
\item proporzionalità diretta
\item crescita lineare
\item nessuna delle altre risposte
}

\itemdomanda{Considerando i dati in tabella l'unica relazione possibile tra le grandezze $x$ e $y$ è di:
\renewcommand{\arraystretch}{1.5}
\begin{tabular}{|>{\columncolor{headercolor}}c|c|c|c|c|c|c|}
\hline
$x$ & 2 & 4 & 6 & 8 & 10 & 12 \\ \hline
$y$ & 101 & 202 & 303 & 404 & 505 & 606 \\ \hline
\end{tabular}

}{
\item proporzionalità quadratica
\item proporzionalità inversa
\item proporzionalità diretta
\item crescita lineare
\item proporzionalità cubica
}

\itemdomanda{Considerando i dati in tabella l'unica relazione possibile tra le grandezze $x$ e $y$ è di:
\renewcommand{\arraystretch}{1.5}
\begin{tabular}{|>{\columncolor{headercolor}}c|c|c|c|c|c|c|}
\hline
$x$ & 2 & 4 & 6 & 8 & 10 & 12 \\ \hline
$y$ & 50 & 25 & $\frac{50}{3}$ & $\frac{25}{2}$ & 10 & $\frac{25}{3}$ \\ \hline
\end{tabular}

}{
\item proporzionalità quadratica
\item proporzionalità inversa
\item proporzionalità diretta
\item crescita lineare
\item nessuna delle altre risposte
}

\itemdomanda{Considerando i dati in tabella l'unica relazione possibile tra le grandezze $x$ e $y$ è di:
\renewcommand{\arraystretch}{1.5}
\begin{tabular}{|>{\columncolor{headercolor}}c|c|c|c|c|c|c|}
\hline
$x$ & 2 & 4 & 6 & 8 & 10 & 12 \\ \hline
$y$ & 4 & 16 & 36 & 64 & 100 & 144 \\ \hline
\end{tabular}

}{
\item proporzionalità quadratica
\item proporzionalità inversa
\item proporzionalità diretta
\item crescita lineare
\item proporzionalità cubica
}

\newpage
\section*{Appendice: risposte corrette}
\begin{longtable}{@{}r>{\raggedright\arraybackslash}p{0.14\linewidth}>{\raggedright\arraybackslash}p{0.20\linewidth}>{\raggedright\arraybackslash}p{0.42\linewidth}c@{}}
\toprule
\textbf{N.} & \textbf{Materia} & \textbf{Argomento} & \textbf{Titolo} & \textbf{Risp.} \\
\midrule
\endfirsthead
\toprule
\textbf{N.} & \textbf{Materia} & \textbf{Argomento} & \textbf{Titolo} & \textbf{Risp.} \\
\midrule
\endhead
\midrule
\multicolumn{5}{r}{\footnotesize\textit{Continua}} \\
\endfoot
\bottomrule
\endlastfoot
1 & Logica & Logica Deduttiva & Relazione transitiva & C \\
2 & Logica & Logica Deduttiva & Sillogismo dei canarini & A \\
3 & Matematica & Tipi di Relazione & Crescita lineare & D \\
4 & Matematica & Tipi di Relazione & Proporzionalità diretta & C \\
5 & Matematica & Tipi di Relazione & Proporzionalità inversa & B \\
6 & Matematica & Tipi di Relazione & Proporzionalità quadratica & A
\end{longtable}

\end{document}
